\subsection{Fusion-Stage Validator Ablation}
To isolate each validator, we re-ran only the fuser on the 30 products with one validator channel removed (\texttt{no\_completeness}, \texttt{no\_accuracy}, \texttt{no\_redundancy}; 90 runs), holding the atomic claims, validator outputs, image, and DOM fixed. The reference here is full \sys, not \base. Table~\ref{tab:ablation} summarizes the structural outcomes.

\begin{table}[t]
\caption{Fusion-stage ablation: primary structural outcomes ($n=30$).}
\label{tab:ablation}
\centering
\footnotesize
\setlength{\tabcolsep}{3pt}
\begin{tabular}{@{}llr@{}}
\toprule
Validator removed & Primary outcome & Value \\
\midrule
Completeness & Validator-added claims lost & 48/48 (100\%) \\
Completeness & Mean $\Delta$ claims / words & $-$1.37 / $-$9.47 \\
Accuracy     & Products with corrections/deletions & 7/30 \\
Redundancy   & Mean $\Delta$ avoidable redundancy rate & +0.048 \\
\bottomrule
\end{tabular}
\end{table}

\textit{Completeness} has the broadest effect. It added at least one claim for 27 of 30 products (48 claims in total, 1.60 per product). Without it, all 48 claims disappear, and outputs lose 1.37 claims and 9.47 words on average; for example, ``short sleeves and a crew neckline'' drop out of a T-shirt description. \textit{Accuracy} interventions are sparse but decisive: 8 corrections and 1 deletion across 7 products. For instance, it corrects ``black contrasting collar'' to ``gray collar with a black inner collar'' and removes the unsupported claim ``a sofa bed is arranged as a single bed''. Without the validator, these errors return. \textit{Redundancy} has a moderate effect: without it, more avoidable DOM repetition is retained in 14 of 30 products. Ablated variants score higher on post-hoc efficiency because validators deliberately add coverage and corrected wording. This trade-off is expected and does not mean that removing validators improves \sys.
